\section{注意力机制}

注意力（Attention）是近年来神经网络领域最重要的创新之一。Manning 教授强调，前面课程中介绍的所有技术（前馈网络、RNN、LSTM、CNN）都是上世纪发明的，而注意力是2014年在神经机器翻译背景下\textbf{全新发明}的概念，它对提升神经网络的能力具有变革性意义。

\subsection{动机：序列到序列的信息瓶颈}

\begin{figure}[H]
\centering
\includegraphics[width=0.95\textwidth]{slides-images/slide-006.jpg}
\caption{Seq2Seq 的信息瓶颈：Encoder RNN 将法语源句 ``il a m' entart\'{e}'' 编码为一个固定向量（橙色框），Decoder RNN 必须仅凭这个向量生成完整的英语翻译 ``he hit me with a pie''\protect\footnotemark}
\end{figure}
\footnotetext{Slides 第7页。}

\begin{figure}[H]
\centering
\includegraphics[width=0.95\textwidth]{slides-images/slide-007.jpg}
\caption{信息瓶颈问题的明确表述：编码器的最后一个隐藏状态必须捕获源句的\textbf{全部}信息，这构成了信息瓶颈\protect\footnotemark}
\end{figure}
\footnotetext{Slides 第8页。}

标准 Seq2Seq 模型的核心问题在于：

\begin{itemize}
    \item 编码器将整个源句压缩为\textbf{单一向量}
    \item 对于长句子，这个向量无法承载所有必要信息
    \item 粗暴的解决方案（增大隐藏状态、增加层数）效果有限
    \item 人类翻译时会\textbf{回头查看}源句的不同部分，而不是一次性记住所有内容
\end{itemize}

\begin{importantbox}{注意力的核心思想}
在解码器的\textbf{每一步}，建立与编码器的\textbf{直接连接}，让解码器能够\textbf{聚焦}于源序列的特定部分，从而按需获取信息。这就像人类翻译时回头查看原文的特定词语一样。
\end{importantbox}

\begin{figure}[H]
\centering
\includegraphics[width=0.95\textwidth]{slides-images/slide-008.jpg}
\caption{注意力机制的核心思想：在解码器的每一步，使用与编码器的直接连接来聚焦于源序列的特定部分\protect\footnotemark}
\end{figure}
\footnotetext{Slides 第9页。}

\subsection{注意力机制的工作流程}

注意力机制在 Seq2Seq 中的完整工作流程可以分解为以下步骤：

\subsubsection{Step 1: 计算注意力分数}

在解码器的每个时间步 $t$，用解码器当前隐藏状态 $\mathbf{s}_t$ 与编码器每个位置的隐藏状态 $\mathbf{h}_i$ 计算注意力分数（最简单的方式是点积）。

\begin{figure}[H]
\centering
\includegraphics[width=0.95\textwidth]{slides-images/slide-009.jpg}
\caption{Step 1：解码器隐藏状态（绿色）与编码器各位置的隐藏状态（红色）做点积，得到注意力分数（蓝色圆点）\protect\footnotemark}
\end{figure}
\footnotetext{Slides 第10页。}

\begin{figure}[H]
\centering
\includegraphics[width=0.95\textwidth]{slides-images/slide-012.jpg}
\caption{继续计算：对编码器的每个位置都计算一个注意力分数\protect\footnotemark}
\end{figure}
\footnotetext{Slides 第13页。}

\subsubsection{Step 2: 通过 Softmax 获得注意力分布}

将注意力分数通过 softmax 转换为概率分布，表示解码器在当前时间步应该``关注''源句哪些位置。

\begin{figure}[H]
\centering
\includegraphics[width=0.95\textwidth]{slides-images/slide-013.jpg}
\caption{Step 2：注意力分数通过 softmax 变为注意力分布（概率分布）。在翻译第一个词时，模型主要聚焦于编码器第一个隐藏状态（对应``il''即``he''）\protect\footnotemark}
\end{figure}
\footnotetext{Slides 第14页。}

\subsubsection{Step 3: 计算注意力输出}

用注意力分布对编码器隐藏状态做\textbf{加权求和}，得到注意力输出（attention output），也称为上下文向量（context vector）。

\begin{figure}[H]
\centering
\includegraphics[width=0.95\textwidth]{slides-images/slide-014.jpg}
\caption{Step 3：用注意力分布对编码器隐藏状态做加权求和，得到注意力输出。该输出主要包含获得高注意力权重的隐藏状态的信息\protect\footnotemark}
\end{figure}
\footnotetext{Slides 第15页。}

\subsubsection{Step 4: 拼接并生成输出词}

将注意力输出与解码器隐藏状态\textbf{拼接}（concatenate），送入后续层生成下一个词。

\begin{figure}[H]
\centering
\includegraphics[width=0.95\textwidth]{slides-images/slide-015.jpg}
\caption{Step 4：将注意力输出与解码器隐藏状态拼接，计算 $\hat{y}_1$，通过 softmax 生成第一个翻译词 ``he''\protect\footnotemark}
\end{figure}
\footnotetext{Slides 第16页。}

\subsubsection{重复：逐步翻译}

在后续的每个解码步骤中，重复上述过程。每一步解码器都会重新计算注意力，关注源句的不同部分。

\begin{figure}[H]
\centering
\includegraphics[width=0.95\textwidth]{slides-images/slide-016.jpg}
\caption{第二个时间步：解码器关注编码器的不同位置，注意力分布发生变化（此时主要关注 ``entart\'{e}''，即 ``pie'' 的动词形式），生成 ``hit''\protect\footnotemark}
\end{figure}
\footnotetext{Slides 第17页。}

\begin{figure}[H]
\centering
\includegraphics[width=0.95\textwidth]{slides-images/slide-017.jpg}
\caption{第三个时间步：注意力分布转移到 ``m'''（法语的 ``me''），生成 ``me''\protect\footnotemark}
\end{figure}
\footnotetext{Slides 第18页。}

\begin{figure}[H]
\centering
\includegraphics[width=0.95\textwidth]{slides-images/slide-020.jpg}
\caption{第六个时间步：翻译 ``pie'' 时，注意力仍然聚焦在 ``entart\'{e}''（法语的 pie/to pie 的动词形式）\protect\footnotemark}
\end{figure}
\footnotetext{Slides 第21页。}

\subsection{本章小结}

注意力机制的工作流程可以总结为四步：
\begin{enumerate}
    \item \textbf{计算注意力分数}：解码器隐藏状态与编码器各位置隐藏状态的相似度
    \item \textbf{Softmax}：将分数转换为概率分布
    \item \textbf{加权求和}：用注意力分布对编码器隐藏状态求加权平均，得到上下文向量
    \item \textbf{拼接与生成}：将上下文向量与解码器隐藏状态拼接，生成输出词
\end{enumerate}

%% ========== Part 3: 注意力的数学形式化 ==========

